% ============================================
% MEM-24
% Gerado de molambudos.md — build_miolo.py
% ============================================
\section*{\textbf{MEM-24} --- A Enfermeira que Chorava}
\markboth{MEM-24 --- A Enfermeira que Chorava}{}

\noindent \textit{Pequenas bondades}

\noindent \textbf{1964 --- Pequenas bondades}

\vspace{0.5em}\noindent\rule{\textwidth}{0.5pt}\vspace{0.5em}

\noindent De vez em quando, uma enfermeira nova aparecia.

\noindent Vinham do interior, moças de famílias pobres que precisavam de trabalho. Entravam no Colônia sem saber o que era. Saíam em uma semana --- ou ficavam e endureciam.

\noindent Uma delas, em 1964, era diferente.

\noindent Ela se chamava Isabel. Tinha 19 anos. Cabelo preto, olhos grandes, mãos que tremiam quando entregavam a comida.

\noindent Na primeira semana, ela passou na frente da minha cela e viu o diário.

\noindent \textit{\textquotedblleft{}O que o senhor escreve?\textquotedblright{}}

\noindent \textit{\textquotedblleft{}A verdade.\textquotedblright{}}

\noindent Ela não entendeu. Mas continuou passando.

\noindent Toda tarde, ela trazia a comida. Enquanto os outros enfermeiros jogavam o prato pelo buraco da porta, ela abria a porta e entregava na mão.

\noindent \textit{\textquotedblleft{}O senhor não é perigoso\textquotedblright{}}, ela disse um dia.

\noindent \textit{\textquotedblleft{}Como sabe?\textquotedblright{}}

\noindent \textit{\textquotedblleft{}Pessoas perigosas não escrevem. Pessoas perigosas não pensam. O senhor escreve e pensa. O senhor é só triste.\textquotedblright{}}

\noindent \textit{\textquotedblleft{}Triste?\textquotedblright{}}

\noindent Ela abaixou os olhos.

\noindent \textit{\textquotedblleft{}O senhor tem o olho mais triste que eu já vi.\textquotedblright{}}

\noindent \textit{\textquotedblleft{}É amarelo.\textquotedblright{}}

\noindent Ela sacudiu a cabeça.

\noindent \textit{\textquotedblleft{}Amarelo é a cor. Triste é o que está por trás.\textquotedblright{}}

\noindent Ela saiu.

\noindent Naquela noite, a criatura falou:

\noindent \textit{\textquotedblleft{}Ela gosta de você.\textquotedblright{}}

\noindent \textit{\textquotedblleft{}Cala a boca.\textquotedblright{}}

\noindent \textit{\textquotedblleft{}Ela quer te salvar.\textquotedblright{}}

\noindent \textit{\textquotedblleft{}Ninguém salva ninguém aqui.\textquotedblright{}}

\noindent \textit{\textquotedblleft{}Ela quer tentar.\textquotedblright{}}

\noindent \textit{\textquotedblleft{}...Ela vai embora.\textquotedblright{}}

\noindent \textit{\textquotedblleft{}Como sabe?\textquotedblright{}}

\noindent \textit{\textquotedblleft{}Todo mundo vai. Todo mundo que entra aqui, sai. Ou de pé. Ou de caixão.\textquotedblright{}}

\vspace{0.5em}\noindent\rule{\textwidth}{0.5pt}\vspace{0.5em}

\noindent Isabel ficou três meses.

\noindent Um dia, ela não veio mais.

\noindent O enfermeiro que trouxe a comida disse: \textit{\textquotedblleft{}Pediu demissão. Disse que não aguentava ver.\textquotedblright{}}

\noindent Perguntei: \textit{\textquotedblleft{}Ver o quê?\textquotedblright{}}

\noindent Ele me olhou.

\noindent \textit{\textquotedblleft{}Ver vocês.\textquotedblright{}}

\noindent Depois que Isabel foi embora, eu fechei o diário por uma semana.

\noindent Não por falta de palavras.
Por falta de esperança.

\noindent Se até quem queria ajudar ia embora, então não adiantava nada.

\noindent O Colônia não era um hospital.
Era um buraco.
E tudo que caía no buraco --- esquecia.

\vspace{0.5em}
\noindent\textit{\textbf{↪ Rotas:} \rota{MEM-25} \rotasep \rota{MEM-12}}